% New independently checked contribution after checkpoint05.
% Root assigns the final A.4 statement numbers; labels and NJ tags are stable.
\subsection{The exact first-order normalizer with the signed period retained}

\begin{result}{A.4.20}{Proved}{Both determinant and spinor can be corrected inside the centralizer}
\label{lem:orthogonal-centralizer-spinor}
For the actual subgroup \(H\simeq C_3^2\) on the eighteen-dimensional
quaternion multiplicity space \((M,B)\), the map
\[
 C_{O(M,B)}(H)\longrightarrow
 \{1,-1\}\times\mathbb F_3^\times/(\mathbb F_3^\times)^2,
 \qquad c\longmapsto(\det c,\theta(c))
 \tag{NJ.1}
\]
is surjective. In particular its special orthogonal centralizer contains
an element of nonsquare spinor norm. Right multiplication by these
centralizer elements adjusts both invariants of an orthogonal conjugator
while preserving its exact induced map on \(H\).
\end{result}
\begin{proof}
Use the already proved free module \(M=A^2\), where
\(A=\mathbb F_3[H]\), the involution \(\bar g=g^{-1}\), and the
Hermitian form
\[
 h(x,y)=\sum_{g\in H}B(x,gy)g.
\]
Write \(\epsilon:A\to\mathbb F_3\) for augmentation. For a vector
\(v\) with unit norm \(a=h(v,v)\), the actual Hermitian reflection
\[
 R_v(x)=x-2h(x,v)a^{-1}v
 \tag{NJ.2}
\]
centralizes \(H\), preserves \(B\), and is minus one on the
nine-dimensional nondegenerate summand \(Av\), with identity on its
orthogonal complement. Thus \(\det R_v=-1\).

In the basis \(\{gv:g\in H\}\) of \(Av\), its bilinear Gram matrix
is the regular multiplication matrix of the self-adjoint element \(a\).
Multiplication by \(a\) acts by \(\epsilon(a)\) on every graded piece
of the augmentation filtration. Their dimensions sum to nine, so
\[
 \det(B|_{Av})=\epsilon(a)^9=\epsilon(a).
\]
Diagonalizing this nondegenerate form expresses \(R_v\) as the
product of the nine ordinary orthogonal reflections in an orthogonal
basis of \(Av\). The quadratic form is \(q(x)=\tfrac12 B(x,x)\).
The reflection formula for spinor norm, and its multiplicativity,
therefore give
\[
          \theta(R_v)=2^9\det(B|_{Av})=2\epsilon(a).
 \tag{NJ.3}
\]
These are the definition and Theorem11.50 on printed pages163--164
of \cite{GhostTaylorClassical}; the formula uses this specified
quadratic form throughout.

The reduced Hermitian form on \(M/\mathfrak mM\) is a nondegenerate
symmetric binary form. It represents both nonzero elements of
\(\mathbb F_3\): diagonalize it as \(aX^2+bY^2\); unequal
coefficients already give both values, and equal coefficients give
the two values at \((1,0)\) and \((1,1)\).
Lift vectors of norms one and two to \(M\). Formula (NJ.3) gives
two centralizer elements whose determinant/spinor pairs are
\((-1,2)\) and \((-1,1)\). Their product has pair \((1,2)\).
Together they give all four values of (NJ.1).
The identity
\((gc)h(gc)^{-1}=ghg^{-1}\) for \(c\in C_{O(M)}(H)\)
proves the final assertion with the induced map retained.

The independent program \texttt{check\_main\_spinor\_correction.py}
constructs these two Hermitian reflections in the actual eighteen
coordinates and verifies their product again by an ordinary-reflection
factorization of length twenty-four. Its saved matrices satisfy
\(c^tBc=B\), \(ch_i=h_ic\), \(\det c=1\), and \(\theta(c)=2\).
\end{proof}

Let \(s,t\) be the two specified translation generators of \(H\), and put
\[
 \begin{gathered}
 A=\mathbb F_3[X,Y]/(X^3,Y^3),\qquad X=s-1,\quad Y=t-1,\\
 \mathfrak m=(X,Y),\qquad
 M_1=M/\mathfrak m^2M,\qquad
 A_1=A/\mathfrak m^2=\mathbb F_3[X,Y]/(X,Y)^2.
 \end{gathered}
 \tag{NJ.4}
\]
The actual quotient \(M\to M_1\) is the multiplicity map of
the compact faithful detector: its kernel is exactly
\(\mathfrak m^2M\), as proved in the preceding detector result.
In particular \(\dim M_1=6\). This is an induced map on the
specified source module.

Write
\[
 P=\{q\in\mathrm{SL}_2(\mathbb Z/12):q\equiv I\pmod2\}.
 \tag{NJ.5}
\]
For its mod-four coordinate use
\[
 q\equiv I+2\begin{pmatrix}\alpha&\beta\\\gamma&\alpha\end{pmatrix}
 \pmod4,\qquad
 \eta(q)=(-1)^{\alpha+\beta+\gamma}.
 \tag{NJ.6}
\]
Multiplication adds the three displayed binary coordinates, so
\(\eta\) is a character. The mod-three and mod-four coordinates
are independent by the Chinese remainder theorem, and consequently
\(|P|=24\cdot8=192\).
These matrices use the downstairs handle basis \((a,b)\).
The primitive upstairs handle basis is \((b^3,a^3)\); the exact
change of basis is the swap \(J=\left(\begin{smallmatrix}0&1\\1&0
\end{smallmatrix}\right)\), so \(q_{\rm down}=Jq_{\rm up}J\).
The character (NJ.6) is invariant under this swap.

For \(g\in N_{\mathrm{SO}(M,B)}(H)\), let
\(\nu_g\in\mathrm{GL}_2(\mathbb F_3)\) be its actual action on
the labelled generators \((s,t)\), and let \(\bar g\) be its
induced action on \(M_1\). Define the finite group
\[
 \mathcal N_1=
 \{(\bar g,q):g\in N_{\mathrm{SO}(M,B)}(H),\ q\in P,
       \ \nu_g=q\bmod3,\ \theta(g)=\eta(q)\}.
 \tag{NJ.7}
\]
The signed period is part of the displayed pair.

\begin{result}{A.4.21}{Proved}{The full allowed first-order image is the ordinary image}
\label{prop:orthogonal-first-order-joint-normalizer}
The group \(\mathcal N_1\) in (NJ.7) has order
\[
                       |\mathcal N_1|=384\cdot3^6=279936.
 \tag{NJ.8}
\]
It is exactly the image on \(M_1\), together with signed period
modulo twelve, of the kernel of the quaternion automorphism action
in the group generated by the four actual substitutions
\[
 \begin{array}{c|cc}
       &a&b\\ \hline
 u&a&ba\\
 l&ab^{-1}&b\\
 \operatorname{Inn}(a)&a&aba^{-1}\\
 \operatorname{Inn}(b)&bab^{-1}&b.
 \end{array}
 \tag{NJ.9}
\]
Thus every element of (NJ.7) has an ordinary modular/deck word with
exactly the same action on this six-dimensional quotient and exactly
the same signed period modulo twelve.
\end{result}
\begin{proof}
We first prove the upper bound intrinsically. Let \(H_0\) be the
augmentation of the Hermitian Gram matrix in a free \(A\)-basis of
\(M\). In the group basis of \(A\), the coefficient-of-identity
pairing \((a,b)\mapsto[a\bar b]_1\) is the identity matrix.
Consequently the determinant of the bilinear form underlying the
Hermitian Gram matrix is its determinant as an \(\mathbb F_3\)-linear
operator. On each augmentation-graded piece this operator is
\(H_0\), repeated according to that piece's dimension. It follows that
\[
                       \det B=(\det H_0)^9.
 \tag{NJ.10}
\]
Change of an \(\mathbb F_3\)-basis multiplies this determinant by a
nonzero square, which is one. The actual \(B\) has determinant one,
so \(\det H_0=1\). The binary form \(H_0\) is therefore anisotropic:
an isotropic nondegenerate binary form has determinant minus a square,
which here is two. In the explicitly checked free basis obtained from
standard columns one and two, counted from zero,
\(H_0=\operatorname{diag}(2,2)\).

An \(H\)-normalizer \(g\) is \(\nu_g\)-semilinear over \(A\), and
the definition of the Hermitian form gives the exact identity
\[
                  h(gx,gy)=\nu_g(h(x,y)).
 \tag{NJ.11}
\]
Its constant term \(T_0\) therefore preserves \(H_0\).
Write \(g\) as an invertible \(A\)-matrix followed by the ring
permutation induced by \(\nu_g\) on each of the two free summands.
The latter has square determinant over \(\mathbb F_3\).
The former, by its augmentation filtration, has determinant
\((\det T_0)^9\). Thus
\[
             \det_{\mathbb F_3}g=(\det T_0)^9,
 \qquad T_0\in\mathrm{SO}(H_0)\simeq C_4.
 \tag{NJ.12}
\]
The last group can also be read directly from
\(H_0=\operatorname{diag}(2,2)\): its four matrices are
\(I,-I,\left(\begin{smallmatrix}0&1\\-1&0\end{smallmatrix}\right)\)
and the negative of the latter.

For fixed \((\nu_g,T_0)\), differences of two allowed first-order
actions are \(A_1\)-linear, with the form
\(I+XD+YE\), where \(D,E\in M_2(\mathbb F_3)\).
Here \(\bar X=-X\) and \(\bar Y=-Y\) in \(A_1\).
Keeping all first-degree terms in the exact isometry equation gives
\[
                  D^tH_0=H_0D,\qquad E^tH_0=H_0E.
 \tag{NJ.13}
\]
For each coefficient the map \(D\mapsto H_0D\) identifies this
space with the symmetric binary matrices, of dimension three.
There are at most \(3^6\) first-order actions for fixed
\((\nu_g,T_0)\).

The kernel of constant reduction in the \(A\)-linear isometry group
is contained in \(1+M_2(\mathfrak m)\), a finite group of order a
power of three. Indeed its filtration by
\(1+M_2(\mathfrak m^i)\) has elementary abelian successive
quotients. Spinor norm is therefore trivial on this full kernel.
Reduction of the special orthogonal centralizer onto \(\mathrm{SO}(H_0)\)
is surjective: each ordinary reflection of \(H_0\) lifts by (NJ.2),
and products of two such lifts give its four rotations.
Lemma~\ref{lem:orthogonal-centralizer-spinor} shows that spinor norm
induces a nontrivial character on this \(C_4\).
Every nonempty fibre with a fixed \(\nu_g\) is a coset of this
centralizer. Hence, for each fixed \(q\in P\), the prescribed
spinor in (NJ.7) allows exactly two possible constant terms among
those of that fibre. Each has at most \(3^6\) first-order lifts.
This proves
\[
                       |\mathcal N_1|\le192\cdot2\cdot3^6.
 \tag{NJ.14}
\]

For the reverse count we use the complete finite certificate with
the source maps specified. The substitutions (NJ.9) have quaternion
automorphism image of order twenty-four. Enumerating its cosets
and taking every coset-generator edge gives ninety-six actual
Schreier words. Their action on the rank-seventy-three free cover
is obtained by substitution followed by literal Schreier rewriting.
For each word the retained program restricts that matrix to the
eighteen-dimensional multiplicity space, and directly verifies
\[
 \begin{gathered}
        g^tBg=B,\qquad\det g=1,\qquad
        gh_i g^{-1}=h_1^{q_{1i}}h_2^{q_{2i}},\\
        \theta(g)=\eta(q),\qquad
        \Phi g=\bar g\Phi.
 \end{gathered}
 \tag{NJ.15}
\]
Here \(\Phi:M\to M_1\) is the actual source-word map of the
compact detector, of rank six and kernel \(\mathfrak m^2M\).
The period \(q\) is obtained by multiplying the two ordinary
period matrices of \(u,l\), with identity for the two inner
substitutions. These checks give inclusion of the whole generated
joint image in (NJ.7), with the two coordinates matched word by word.

For completeness the enumeration uses a split linear coordinate
system, so it does not enumerate an unspecified ambient group.
Choose two free generators of \(M_1\) over \(A_1\), and order its
basis as these generators, their \(X\)-multiples, then their
\(Y\)-multiples. For \(T_0\in\mathrm{GL}_2(\mathbb F_3)\) and
\(\nu\in\mathrm{GL}_2(\mathbb F_3)\), the split semilinear matrix is
\[
 S(T_0,\nu)=\operatorname{diag}(T_0,\nu\otimes T_0).
 \tag{NJ.16}
\]
Every source matrix has a unique expression
\((I+L_C)S(T_0,q\bmod3)\), where the only nonzero block of
\(L_C\) is its lower-left \(4\)-by-\(2\) block \(C\).
Thus the additive kernel is a subspace of \(\mathbb F_3^8\),
with the exact base action
\[
 C\longmapsto(\nu\otimes T_0)C T_0^{-1}.
 \tag{NJ.17}
\]
Breadth-first multiplication of the ninety-six base generators
\((T_0,q)\) gives exactly \(384\) distinct values. With one actual
source representative for each, all \(384\cdot96=36864\) Schreier
edges give vectors in the additive kernel. Row reduction gives
rank six. The same program checks that this six-dimensional span
is stable under every generator by (NJ.17). Schreier's rewriting
identity proves that these vectors generate the entire kernel;
stability proves closure of the displayed extension. The ordinary
joint image therefore has order \(384\cdot3^6\).
Together with (NJ.14) and the inclusion from (NJ.15), this proves
equality and every assertion of the proposition.

The independent verifier
\texttt{check\_main\_normalizer\_joint.py} rebuilds the four
rank-seventy-three maps from the original free words, forms all
ninety-six words itself, and factors their spinor norms into ordinary
reflections. It checks all intertwiners against the independent
source-to-detector map before performing the \(384\)-state and
\(36864\)-edge calculations. Its actual output, input hashes,
matrices, reduced binary form and first-order kernel are saved.
The ninety-six spinor values consist of seventy-two squares and
twenty-four nonsquares. All six computed kernel directions satisfy
(NJ.13).
\end{proof}
